% Figure: free-body diagram of a car on a banked curve. Normal force,
% weight, friction, and the required centripetal direction.
\begin{figure}[htbp]
\centering
\begin{tikzpicture}[
  force/.style={draw=engblue, -{Stealth}, very thick},
  wt/.style={draw=engblack, -{Stealth}, very thick},
  fr/.style={draw=engamber, -{Stealth}, very thick},
  cen/.style={draw=engred, -{Stealth}, thick, dashed},
  lbl/.style={font=\scriptsize},
]
% bank angle 18 degrees. Road surface line.
\def\bank{18}
% ground / banked road
\draw[draw=enggray, thick] (-3,0) -- (3,1.949);  % tan(18)=0.3249, slope over 6 -> 1.949
\fill[enggray!12] (-3,0) -- (3,1.949) -- (3,-0.4) -- (-3,-0.4) -- cycle;
% car block sitting on the bank near centre, rotated
\begin{scope}[shift={(0,0.6494)}, rotate=\bank]
  \draw[draw=engblack, thick, fill=engblue!10, rounded corners=1.5pt]
    (-0.7,0) rectangle (0.7,0.45);
\end{scope}
% point of application = centre of car base on the road
\coordinate (P) at (0,0.6494);
% Normal force: perpendicular to road surface (road slopes up-right at 18 deg,
% so normal points up-left at 18 deg from vertical): direction (−sin18, cos18)
\draw[force] (P) -- ++({-1.6*sin(\bank)}, {1.6*cos(\bank)})
  node[above, lbl, engblue] {$N$};
% Weight straight down
\draw[wt] (P) -- ++(0,-1.6) node[below, lbl] {$mg$};
% Friction up the slope (toward the high side) along surface: (cos18, sin18)
\draw[fr] (P) -- ++({1.3*cos(\bank)}, {1.3*sin(\bank)})
  node[right, lbl, engamber] {$f$ (up-slope)};
% required centripetal direction: horizontal, toward centre of curve (left)
\draw[cen] (P) -- ++(-1.8,0) node[left, lbl, engred] {$a_c=v^2/r$ (toward centre)};
% bank angle arc
\draw[draw=enggray] (1.2,0.78) arc (\bank:0:1.2);
\node[lbl, enggray] at (1.5,0.45) {$\theta$};
\end{tikzpicture}
\caption[Free-body diagram of a car on a banked curve]{Free-body
diagram of a car cornering on a road banked at angle $\theta$. Three
forces act: the weight $mg$ straight down, the normal force $N$
perpendicular to the road surface, and the tyre friction $f$ in the
plane of the road. The acceleration is horizontal and points toward
the centre of the curve with magnitude $v^{2}/r$ (red dashed). Banking
tilts the normal force so that its horizontal component supplies part
of the centripetal requirement without help from friction; at the
design speed $v_{0}=\sqrt{g r\tan\theta}$ the friction is zero and the
normal force alone turns the car. Above that speed friction must point
down-slope to add inward force; below it friction must point up-slope
to keep the car from sliding down the bank. The friction-limited
maximum speed adds $\mu N$ on the inward side and is worked in the text.}
\label{fig:vol03ch04:banked-curve}
\end{figure}
